\documentclass[UTF8,11pt,a4paper]{ctexart}
\usepackage[margin=1.9cm,headheight=15pt]{geometry}
\usepackage{amsmath,booktabs,longtable,array,graphicx,tikz,xcolor,hyperref,xurl,fancyhdr,enumitem}
\usetikzlibrary{arrows.meta}
\setmainfont{Arial}\setmonofont{Menlo}
\setCJKmainfont{Songti SC}\setCJKsansfont{Heiti SC}
\definecolor{ink}{HTML}{163C50}\definecolor{accent}{HTML}{007C91}
\hypersetup{colorlinks=true,urlcolor=accent,linkcolor=ink,pdftitle={S81 真实深度评分增补},pdfauthor={Yiyang Liu -- AI辅助整理}}
\renewcommand{\UrlFont}{\fontspec{Arial Unicode MS}\small}
\setlength{\parskip}{.38em}\linespread{1.05}\setlength{\emergencystretch}{2em}
\setlist{nosep,leftmargin=1.7em}\renewcommand{\arraystretch}{1.12}
\pagestyle{fancy}\fancyhf{}\lhead{\small CSIT 6910：S81 真实深度评分增补}\rhead{\small 2026-09-11 01:11 实际结果}\cfoot{\small\thepage}
\ctexset{section={format=\large\bfseries\sffamily\color{ink},beforeskip=.8em,afterskip=.4em}}
\begin{document}
\begin{center}{\LARGE\bfseries 从“一条可能位置线”\\到“一个几何期望点”}\par\vspace{.3em}{\large S81：加入真实源深度后的独立增补}\end{center}
\noindent\textbf{先说结果。}这次新增了传感器深度评分，没有生成新图、运行新匹配器或新神经网络。源图1313个特征中，968个可以取到有效深度；全部24行旧匹配保留。在生成图的3828条有效投影记录中，$\le10$像素的计数为0；另1282条生成匹配因深度为0无法评分，不能一并判错。实拍对照也存在尾部异常。

\noindent\textbf{核验状态：}已完成计算、团队内不同作者数值复算及本增补内容与版面检查。

\section{ 新手先抓住：线允许很多位置，深度再选一个位置}
相机关系告诉我们：源图上的一个点，在目标图中应当落在某条线附近，这条线叫\textbf{极线}。沿着这条线，仍有很多可能的位置。源深度提供这个点离相机有多远，于是在固定相机、静态表面等假设下，我们得到一个\textbf{几何期望点}。它是评分器的计算位置，不是生成模型新画出的点。

\begin{center}
\begin{tikzpicture}[x=1cm,y=1cm,>=Latex,font=\small]
\draw[thick,accent] (0,0)--(12,0) node[right]{极线};
\fill[ink] (3,0) circle(2.4pt) node[below=6pt]{期望点 $(25,40)$};
\draw[red,thick] (8.87,-.13)--(9.13,.13) (8.87,.13)--(9.13,-.13);
\node[below=6pt] at (9,0){候选点 $(55,40)$};
\draw[<->,thick] (3,.55)--(9,.55) node[midway,above]{点到点：30 像素};
\node at (6,-1){两点都在线上，候选点到线仍为0。};
\end{tikzpicture}
\end{center}
\noindent\textbf{人工教学示意，非真实数据。}设焦距100像素、主点$(50,40)$，源点在光轴上且深度2米；目标相机向右移动0.5米。按针孔模型，固定点应左移$100\times0.5/2=25$像素，到$(25,40)$。候选$(55,40)$虽然在线上，仍偏离期望30像素。这说明“线距小”不能替代“点位置对”。

\section{ 实际做了什么}
本机在\textbf{北京时间09-11 01:11:10.232329至01:11:10.414074}执行评分，内部计时0.181755秒；不含Python导入等启动时间。解码1张已有真实深度PNG，复用S80的源特征、全部接受匹配和原请求相机。24行、7757条匹配记录完成，28组共有源点比较保存。\textbf{0新RGB解码、0模型调用、0新生成、0训练。}这是新的深度测量，数据场景和生成图均已见过。7757条记录可能重叠，不是7757个独立物理点，更不是独立场景。
\clearpage
\section{ 全部24行：先看支持，再看误差}
所有行固定源特征总数$N=1313$。$M$是S80原接受数；$V$是深度及投影有效数；out是$V$中期望点出目标视野的数量，\textbf{这些点仍计入主误差}。本次每行的$M-V$全部是源深度为0，无其他无效原因。中位与p95单位均为原生576图像像素，$\le10$列为\textbf{计数}。

\begin{center}\small
\setlength{\tabcolsep}{5pt}
\begin{tabular}{rrlrrrrrr}
\toprule 目标&图片&匹配器&$M$&$V$&out&中位px&p95 px&$\le10$数\\\midrule
20 & real & BF & 345 & 300 & 0 & 1.875 & 6.903 & 291 \\
20 & real & LG & 637 & 538 & 0 & 1.983 & 6.866 & 521 \\
20 & A0 & BF & 252 & 202 & 4 & 40.703 & 64.077 & 0 \\
20 & A0 & LG & 642 & 504 & 4 & 40.674 & 65.630 & 0 \\
20 & B & BF & 281 & 212 & 2 & 43.781 & 64.213 & 0 \\
20 & B & LG & 685 & 529 & 3 & 44.009 & 65.150 & 0 \\
21 & real & BF & 175 & 165 & 0 & 3.746 & 9.099 & 158 \\
21 & real & LG & 533 & 475 & 0 & 3.830 & 8.371 & 462 \\
21 & A0 & BF & 196 & 148 & 7 & 109.921 & 154.278 & 0 \\
21 & A0 & LG & 600 & 463 & 14 & 109.831 & 161.819 & 0 \\
21 & B & BF & 209 & 158 & 6 & 115.562 & 172.079 & 0 \\
21 & B & LG & 610 & 472 & 14 & 118.438 & 173.510 & 0 \\
22 & real & BF & 90 & 89 & 1 & 3.156 & 14.604 & 83 \\
22 & real & LG & 400 & 384 & 0 & 3.072 & 24.922 & 344 \\
22 & A0 & BF & 105 & 67 & 9 & 199.217 & 248.284 & 0 \\
22 & A0 & LG & 428 & 302 & 42 & 202.104 & 308.498 & 0 \\
22 & B & BF & 61 & 47 & 4 & 195.424 & 258.527 & 0 \\
22 & B & LG & 379 & 287 & 39 & 201.795 & 294.976 & 0 \\
23 & real & BF & 86 & 84 & 1 & 8.271 & 13.219 & 62 \\
23 & real & LG & 381 & 370 & 0 & 8.346 & 15.276 & 251 \\
23 & A0 & BF & 42 & 15 & 8 & 314.255 & 324.288 & 0 \\
23 & A0 & LG & 269 & 164 & 84 & 314.004 & 344.143 & 0 \\
23 & B & BF & 25 & 19 & 9 & 297.321 & 314.925 & 0 \\
23 & B & LG & 326 & 239 & 82 & 296.806 & 315.947 & 0 \\
\bottomrule
\end{tabular}\end{center}
\noindent real为实拍对照，A0/B为两组已有模型生成图；BF与LG是原S80的两种匹配规则，本次没有重新运行它们。p95不是95\%置信区间，表示所报有效误差分布的95\%分位，使用线性插值。

\textbf{读一行：}目标23的A0/BF，原有42条接受匹配，仅15条可评分，8条期望点出视野；中位误差314.255像素。27条未知深度不能填成0误差，也不能因为没有评分而从原42条记录中消失。完整精度见随附S81\_全部24行.csv及项目原始NPZ。
\clearpage
\section{ 所有无效和实拍异常都要保留}
源图\textbf{968/1313}个特征有有效深度；其余\textbf{345个都是零深度缺测}。本次源点域外数为0，非零深度投影到四个目标时的有效数均为968，没有相机后方或非有限投影。由于同一源点被不同图对重复使用，这些源点数与下面的匹配记录数是不同分母。

\begin{center}
\begin{tabular}{rrrrr}\toprule
目标 & 有效源投影 & 在视野内 & 出视野 & 遮挡真值\\\midrule
20&968&964&4&未知\\21&968&941&27&未知\\22&968&860&108&未知\\23&968&769&199&未知\\\bottomrule
\end{tabular}\end{center}
全部匹配记录：$7757=6233$条有效$+1524$条零深度无效。6233条有效记录中有333条期望点出视野，已经保留在误差分布中；它们可能重复使用同一源点。真实对照$V=2405$，生成组$V=3828$；生成组另有1282条无效记录。\textbf{只能说3828条可评分生成记录的$\le10$像素计数为0，不能将5110条原生成匹配全部判错。}

\textbf{实拍对照不是“干净的完美标尺”。}目标20/21实拍中位数约1.88--3.83像素；目标22的LG中位数虽仅3.072，p95却达24.922像素，384个有效点中40个超过10像素。目标23的实拍中位数为8.271/8.346像素；LG的370个有效点中119个超过10像素。全部实拍有效记录中共233条超过10像素。真实控制也可能受误配、深度边缘、时间差或成像近似影响，不能隐藏这些尾部。

\section{ 一个真实分母算例：支持比例不是准确率}
目标23实拍LG有$M=381$、$V=370$，其中251条误差不超过10像素：
\[
\frac{251}{370}\approx67.84\%,\qquad
\frac{251}{381}\approx65.88\%,\qquad
\frac{251}{1313}\approx19.12\%.
\]
依次回答“有效投影中有多少”“原接受记录中有多少”“全部源特征中有多少”。三种比例都可以报告，却不能互相替代。要称\textbf{匹配准确率}，还须知道候选像素确实属于同一个物理点，本次没有这份完整真值。11条深度无效和932个未接受源特征都保留为不同形式的未知，不能直接判错。

同一源feature ID可以连接实拍与生成组，或连接BF/LG，得到28组辅助比较。但共有源点不保证目标匹配到同一表面，交集也有选择偏差，不能用交集替换全24行主表。目标23的BF比较中，A0减实拍仅1个有效共有源点，B减实拍仅4个，不能把这样的中位数外推。
\clearpage
\section{ 为什么加了真实深度，仍不等于三维真值}
源特征坐标为$p=(u,v,1)^T$，光轴深度为$Z$，内参为$K$；光学相机的旋转和中心为$R_s,c_s,R_t,c_t$。本次计算
\[
X_s=ZK^{-1}p,  X_t=R_t^T(R_sX_s+c_s-c_t), 
\hat q=\pi(KX_t),  e=\|q-\hat q\|_2.
\]
$Z$是沿相机光轴的深度，\textbf{不是到相机中心的欧氏距离}。深度PNG原码除5000转为米，零值为缺测；fr2已有深度修正不再乘一次。$e$是二维像素误差，不是米、旋转角或相机轨迹误差。

\begin{enumerate}
\item \textbf{不同步仍在。}唯一深度时间为1311868171.680537，RGB为1311868171.663411，深度晚17.126毫秒。本次沿用旧关联、RGB时刻的相机，不作运动补偿；像素已注册不代表同时曝光。
\item \textbf{内参和坐标仍为旧声明。}沿用近似ROS K、插值光学c2w和未额外去畸变的fr2数据；576逆映射为$(u+96)/1.2,v/1.2$。最近邻用$\lfloor x+0.5\rfloor$，不给无效点寻找有利替代深度，也不看评分修K或尺度。
\item \textbf{边界和深度赋值是近似。}原网格连续域限$0\le u\le639,0\le v\le479$，比$[0,640)\times[0,480)$更保守。最近邻Z被赋给未取整特征射线，是局部深度常量近似；边缘混合和RGB面积插值不会因此消失。
\item \textbf{可见性未知。}目标Z为正、预测位置在画内，不保证没有被遮挡。本次没有读取目标深度，不称遮挡感知评分。出视野记录也不能直接解释为已知遮挡。
\end{enumerate}

大误差限制的是“同一静态表面、深度与时间有效、匹配正确、当前成像和请求相机合理”的\textbf{联合假设}，不能唯一定位为相机、3D表示或记忆模块故障；小误差也不证明全图正确。

\section{ 新测量与独立复核的边界}
S80主量是\textbf{双向点到线距离的平均}；这里是\textbf{前向点到点距离}，不能把同一个10像素标签当成同一个测量。使用同一非退化K/相机模型时，期望点应在目标极线上，故$e\ge d_{\rm target-line}$；此下界不适用于S80双向均值。本次源码在执行时检查了目标极线恒等和这条正确下界，数值容差不是物理精度认证。

\noindent\textbf{核验状态：}已完成计算、团队内不同作者数值复算及本增补内容与版面检查。

不同作者已用独立标量实现全量复算源取样、四目标投影、全部匹配、配对与CSV，root已接受条件性结果；这不是外部模型复现或传感器标定。作者按实际回执和完整摘要制作报告，并从保存的SOURCE\_SAMPLES读取原因计数；没有重新运行评分、原图或深度解码。全部原匹配、无效值与完整精度由原实验目录保留。
\clearpage
\section{ 下一步：让普通几何引导真正进入生成}
S81把已见失败量得更具体，但没有改变生成器。下一实现选择是\textbf{普通的历史几何重投影加遮罩生成引导}：用允许历史预测几何，渲染目标视角可见内容；在一段预先固定的去噪阶段，对目标槽的干净预测作空间融合。它是候选实现路线，\textbf{尚未执行，也不是已验证的新方法}。下表的warp指投影后的历史图，mask指有来源支持的区域遮罩。

\begin{center}\small
\begin{tabular}{p{2.0cm}p{7.2cm}p{6.0cm}}\toprule
条件&做什么&比较目的\\\midrule
$G_0$&原A0生成；可证明输入与实际随机流一致才复用。&给出不加引导的基线。\\
$G_{\rm paste}$&同一几何warp在同一mask内与$G_0$作末端像素合成，洞内保留$G_0$。&检验是否仅靠最后贴回历史像素就能解释改善。\\
$G_{\rm guide}$&采样过程中把同一warp内容融合进目标干净预测，再完成普通去噪。&检验生成过程是否带来超出末端复制的改善。\\\bottomrule
\end{tabular}\end{center}

首轮保留A0的四张历史$[19,18,13,12]$、四目标$[20,21,22,23]$、请求相机、模型/VAE和原采样步数。主臂几何只能来自这些\textbf{合法历史RGB}及允许相机；缓存来源、尺度和坐标先核清。S81的传感器深度仍用于评分，不能悄悄加入RGB-only主臂；若另加RGB-D臂，必须承认新增模态。目标RGB/depth和评分不能用来挑warp、调mask、强度或时段。

若$G_{\rm paste}$已经解释全部改善，或$G_{\rm guide}$改善覆盖区却恶化洞、边界或内容，就不能称生成控制取得进步。几何预处理、渲染、编码和采样成本分别记录；不能把额外引导算力当记忆方法收益。完整WorldForge等强近邻仍是更高的方法对照，简化基线胜出也不等于创新成立。

\section{ 已有工作与创新边界}
本轮来源审查确认：WorldForge已有部分图/mask与采样中纠正；Latent-Reframe已有中途重拍与再次去噪；Gen3C已有深度构建3D缓存与渲染条件生成。它们的机制、训练和成本不同，本机简化引导不能冒充完整复现。\textbf{目前仍为NO\_METHOD\_SELECTED，novelty\_authorization=NONE；不保证CCF A或导师评价。}

\noindent\textbf{本轮证据导航：}\href{/Users/rocket/Desktop/HKUST\%20IT/ip-/geometry-world-modeling/work/S81_anchor_depth_reprojection/MEASUREMENT_NEAREST_WORK.md}{测量来源说明}；\href{/Users/rocket/Desktop/HKUST\%20IT/ip-/geometry-world-modeling/work/S81_anchor_depth_reprojection/GENERATION_MECHANISM_NEXT_STEP.md}{下一生成基线路线}；\href{/Users/rocket/Desktop/HKUST\%20IT/ip-/geometry-world-modeling/work/S81_anchor_depth_reprojection/execution_01/ROWS.json}{全部24行原始摘要}；\href{/Users/rocket/Desktop/HKUST\%20IT/ip-/geometry-world-modeling/work/S81_anchor_depth_reprojection/execution_01/ALL_RECORDS.csv}{逐条匹配记录}；\href{/Users/rocket/Desktop/HKUST\%20IT/ip-/geometry-world-modeling/work/S81_anchor_depth_reprojection/ROOT_RESULT_ACCEPTANCE.json}{root科学接受记录}。来源核读范围见同目录两份SOURCE\_SCOPE；原文为不同作者局部核查，本报告未新联网或执行官方代码。

\noindent\small\textbf{一手入口：}
\href{https://cvg.cit.tum.de/data/datasets/rgbd-dataset/file_formats}{TUM深度/位姿约定}；
\href{https://raw.githubusercontent.com/zju3dv/LoFTR/master/src/loftr/utils/geometry.py}{LoFTR几何帮助函数（动态master，未固定commit）}；
\href{https://arxiv.org/html/2509.15130v3}{WorldForge v3}；
\href{https://arxiv.org/html/2412.06029v1}{Latent-Reframe v1}；
\href{https://arxiv.org/html/2503.03751v1}{Gen3C v1}。
\end{document}
